\section{Attachment, curvature signal, and matched observing states}
\label{at:section}

Apply \Cref{thm:dn-comparison} to the actual entry data and the
template from \Cref{prop:packet-preparation}. Write $g_T$ for the
resulting exact finite double-null metric. Its fields differ from
the template by $e^{-(1+c)\alpha T}$ in every preassigned finite
observation order, for a fixed $c>0$. All orders used below,
including thirty ordinary observation derivatives, are included
before choosing $\zeta$. We first place this metric in the full
development of $d_T$.

\subsection{Attaching the actual exterior and interior}

\begin{proposition}[Attachment to the full development]
\label{prop:experiment-attachment}
For sufficiently large finite $T$, the exact exterior and deep
experiments, restricted inside their construction margins, lie
in a smooth globally hyperbolic vacuum development of the entire
datum $d_T$. This development embeds isometrically in its full
MGHD $M_{d_T}$. The full initial bridge and the exterior correction
tail are retained. The common earlier branch tubes and the compact
observation region have interior margins in the attached region.
\end{proposition}

\begin{proof}
The lower and exterior joins are the qualitative construction in
\cite[\CompanionRef{signed}{attach:whole-bridge-section}, \CompanionRef{signed}{attach:geometric-contract}]{KerrCOne}, based on
\cite[\CompanionRef{foundation}{cmp:finite-range}, \CompanionRef{foundation}{cmp:development}, \CompanionRef{foundation}{pk2:attachment-proposition}]{KerrCTwo}. Its hypotheses are full-bridge $p_{10}$
closeness, exact protected agreement on the compact causal
shadows of the early joins, actual corrected partial-slice data,
low exterior closeness on enlarged matching bands, and strict
outflow and overlap margins. We verify these hypotheses and the
new double-null join separately.

The finite lower-block construction in
\cite[\CompanionRef{foundation}{cmp:finite-range}]{KerrCTwo} applies to $d_T$ since
$p_{10}(d_T-d)\to0$. It contains every spatial radius of the new
bridge, a past collar, the two bent exterior initial slices, and
fixed forward bands on them. On the ends the bent slices equal
the original bridge outside a fixed compact set. The protected
ball \eqref{pk:protected} contains a fixed-origin ball of radius
\[
 (1-2\beta-s_*)t_T-O(1)\longrightarrow\infty.
\]
Consequently every fixed compact causal shadow needed for the
bent pieces and bounded connecting pieces is exactly unchanged.
Domain of dependence in a common early Cauchy gauge identifies
these pieces with the background. At large radii the horizontal
selected-end slice carries the actual new datum, including its
tail; no replacement by the old datum is made there.

Choose a fixed exterior joining time $t_a$ inside the enlarged
lower block. On the selected end attach the exact exterior from
\Cref{prop:packet-preparation}, computed from these same
partial-slice data, using the shrinking outer cutoff
\eqref{pk:outer-cutoff}.
The identifying wave map on an enlarged fixed band about $t_a$
has the prescribed initial slice map and matching unit normals.
Fixed-time wave-map stability gives a diffeomorphism on the
smaller band, uniformly through the large outer cutoff: the
exterior charts have uniform low regular geometry, the new
metrics are low-order close, and artificial sides are strictly
outgoing. The induced metric and second fundamental form agree;
the wave gauge then determines the remaining normal metric jets.
Reduced uniqueness therefore makes this map an isometry. These
are exactly the finite-band hypotheses checked in
\cite[\CompanionRef{signed}{attach:whole-bridge-section}, \CompanionRef{signed}{attach:geometric-contract}]{KerrCOne}. The unselected end uses the
unchanged exterior. Straighten and identify inside the matching
margins, without keeping a second copy of overlapping points.

It matters that the lower block has not been spatially truncated.
In a selected-end overlap the retained region has, in addition
to its finite future cuts, the condition
\begin{equation}\label{at:outer-union}
 t_e<t_a\quad\hbox{or}\quad r_{\rm co}<R_{\rm out}(t_e).
\end{equation}
Both $t_e$ and $r_{\rm co}-R_{\rm out}(t_e)$ increase strictly
along future causal curves at their active boundaries. The
derivative of the second quantity with respect to $t_e$ is at least
$V_{\rm out}-\sup|\dd r_{\rm co}/\dd t_e|>0$.
Their strict sublevels and their union are preserved to the past.
A past curve in the added piece stays inside its outer cutoff
until reaching $t_a$, after which all radii are present. A past
curve already below $t_a$ stays there. The omitted corner at
$t_e=t_a$, $r_{\rm co}=R_{\rm out}(t_a)$ is consequently not a
missing past endpoint.

The required interval of $\mathcal S$ and a margin are present
in this joined lower/exterior region. At its selected late end,
$t_e\le t_T+O(1)$ and the co-radius is bounded; its middle and
unselected portions are unchanged. Low closeness preserves its
spacelike character. Use the \emph{actual exterior induced data}
there. The nearby eikonal roots with the prescribed entry values,
and the shift-free angular transport, define the new coordinates
on a collar. They are noncharacteristic at this fixed depth.
Local geometric vacuum uniqueness identifies this collar with
the initial collar of the exact solution from
\Cref{thm:dn-comparison}; equivalently that solution may be
started on the collar itself. The quadratic template comparison
does not substitute for matching actual data. For every fixed
finite $T$ the collar widths can be chosen inside the available
margins; no lower width uniform in $T$ is required.

Above the join all cuts are made in the \emph{new} double-null
coordinates. For a causal vector $V$ in their future cone,
\[
 -2a_T\,\dd u(V)\dd v(V)
 +\gamma_T\bigl(\dd\theta(V)-b_T\dd u(V),
                  \dd\theta(V)-b_T\dd u(V)\bigr)\le0.
\]
Thus $\dd u(V)$ and $\dd v(V)$ have the same weak sign, and
the future choice is the nonnegative sign. Also
$g_T^{-1}(\dd t,\dd t)=-1/(2a_T)<0$; hence $t$ is temporal.
For a nonzero future causal vector, $\dd t(V)>0$, and
\begin{equation}\label{at:deep-exits}
 \frac{\dd}{\dd t}\left(u+\frac{c_0t}{T}\right)
       \ge\frac{c_0}{T},\qquad
 \frac{\dd}{\dd t}\left(v+\frac{c_0t}{T}\right)
       \ge\frac{c_0}{T}.
\end{equation}
The strict side and terminal sublevels of \eqref{pk:domain}
are therefore preserved into the past. In particular a past-inextendible
causal curve from its strict interior reaches $t=0$ strictly inside the
attached interval. Indeed its side functions never exceed their
strictly smaller starting values. On each fixed experiment, the retained
closure lies inside a compact enlargement on which smooth nondegeneracy
bounds $a_T,a_T^{-1},b_T,\gamma_T^{\pm1}$. In $t$ parameter,
$0\le\dd u/\dd t,\dd v/\dd t\le2$, and the displayed cone inequality
then bounds the angular speed using the upper bound for $a_T$ and the
bounds for $b_T,\gamma_T^{-1}$. A finite interior endpoint is therefore
handled by ordinary continuation. This
argument uses the exact new coordinates and does not require
small relative cone error in a deep background wave chart.

Omit every unattached future face of a lower piece, including
its edges. It is a future exit by the spacelike sign, so a past
curve below it cannot end there. At intersections of cuts, the
retained sets are finite unions and intersections of the same
past-preserved strict sublevels; \eqref{at:outer-union} handles
the only change from unrestricted to truncated radii. Matching
collars give product neighborhoods across the included faces.
One first joins the lower bands and then the deep collar at
their intersections, as in \cite[\CompanionRef{foundation}{cmp:development}]{KerrCTwo}; this
gives a Hausdorff smooth manifold with only one copy of each
identified point.

Every past-inextendible causal curve from an added piece therefore reaches
the complete initial bridge. In the lower pieces bounded speeds
exclude escape to spatial infinity in bounded time, and their
Cauchy construction gives the same conclusion. The past collar
gives the future crossing for curves starting below the bridge.
There is at most one crossing, by the lower-block Cauchy property
and the one-way signs of all joins. Thus the complete $\Sigma$
is Cauchy in the assembled vacuum region. Maximal-development
uniqueness embeds this region in $M_{d_T}$, proving the claim.
All restrictions were made inside preassigned enlarged margins,
so the earlier flow tubes and observation neighborhood remain
interior.
\end{proof}

\subsection{A curvature difference on the whole observation patch}

Use the product observation box $\mathcal Q_T$ fixed before the packet
construction in \Cref{pk:section}. It contains all hits of the chosen
seed-neighborhood closure at $v\in T+J$, and the feet of its enlarged
closure lie in the launch-bump plateau where the bump is at least one.
The local polarization there has unit screen norm. These properties
remain available when the seed ball and interval are trimmed below.
Adopt $R(X,Y)=[\nabla_X,\nabla_Y]-\nabla_{[X,Y]}$ and set
\begin{equation}\label{at:D-definition}
 D(g)_{AB}=g\bigl(R(\partial_v,\partial_A)\partial_v,
                                                    \partial_B\bigr).
\end{equation}
Changing the overall curvature convention changes all ensuing
signs together and leaves the norm estimates unchanged.

\begin{proposition}[Whole-patch curvature signal]
\label{prop:curvature-signal}
At every point of $\mathcal Q_T$ with
$|\cos(\lambda v)|\ge1/2$, one has, for all large $T$,
\begin{equation}\label{at:signal}
 |D(g_T)-D(g)|>2\amp.
\end{equation}
The norm is the ordinary angular component norm. No bound on
$D(g)$, or on background parallel-frame curvature, is assumed.
\end{proposition}

\begin{proof}
For a shift-free direction $L_1=\partial_v$, the null connection
identities are
\[
 \nabla_{L_1}L_1=\ell_1L_1,\qquad
 \nabla_{\partial_A}L_1=\chi_{1,A}{}^B\partial_B+
                                                 \xi_{1,A}L_1.
\]
Since $[\partial_v,\partial_A]=0$, applying the curvature
commutator and pairing with $\partial_B$ gives
\[
 D(g)_{AB}=\gamma_{BC}\partial_v(\chi_{1,A}{}^C)
                    +(\chi_1\gamma^{-1}\chi_1)_{AB}
                    -\ell_1\chi_{1,AB}.
\]
Use $\partial_v\gamma=2\chi_1$ to lower the differentiated
index. This gives the exact formula
\begin{equation}\label{at:curvature-formula}
 D(g)=\partial_v\chi_1-\chi_1\gamma^{-1}\chi_1-\ell_1\chi_1
 =\tfrac12\partial_v^2\gamma
   -\tfrac14(\partial_v\gamma)\gamma^{-1}(\partial_v\gamma)
   -\tfrac12(\partial_v\log a)(\partial_v\gamma).
\end{equation}
There is no inverse radial lapse in this formula.

For the template, $\bar\chi_1=\chi_1+\tfrac12\partial_vk_\gamma$
and $\bar\ell_1=\ell_1+\partial_vk_H$ exactly. The linear
variation of \eqref{at:curvature-formula} is therefore
\begin{equation}\label{at:linear-D}
 \begin{split}
 \delta D={}&\tfrac12\partial_v^2k_\gamma
 -\tfrac14\bigl((\partial_vk_\gamma)\gamma^{-1}\partial_v\gamma
 +(\partial_v\gamma)\gamma^{-1}\partial_vk_\gamma\bigr)\\
 &+\tfrac14(\partial_v\gamma)\gamma^{-1}k_\gamma
                              \gamma^{-1}(\partial_v\gamma)
 -\tfrac12\ell_1\partial_vk_\gamma
 -\tfrac12(\partial_vk_H)\partial_v\gamma.
 \end{split}
\end{equation}
The angular Lie term in \eqref{pk:linear-fields} contains only
$Z$ and angular derivatives of $Z$, without a longitudinal
derivative of $Z$. Hence \eqref{pk:coordinate-gain} gives
\[
 \partial_v^2(\Lie_Zg)_{AB}=O(\amp\lambda e^{o(T)}).
\]
The formula for $k_H$ has at most one longitudinal derivative
of $Z$, so $\partial_vk_H=O(\amp\lambda e^{o(T)})$ as well.
All other linear terms in \eqref{at:linear-D} have this bound.
Twice differentiating the real leading oscillation in
\eqref{pk:packet}, whose trace-free leading tensor is unchanged
by trace reversal, yields
\begin{equation}\label{at:leading-signal}
 D(\bar g)-D(g)
 =-\tfrac12\amp\lambda^2\cos(\lambda v)(p_0)_{AB}
      +O(\amp\lambda e^{o(T)})
      +O(\amp^2\lambda^{P}e^{o(T)}).
\end{equation}
The last term comprises quadratic template products in the smooth
relative-field expression \eqref{at:curvature-formula}. Choose
$\zeta$ small enough that this term is $o(\amp)$.
The exact-minus-template error from \Cref{thm:dn-comparison},
at the finite derivative orders in \eqref{at:curvature-formula},
is also $o(\amp)$; multiplication by the individual
$e^{o(T)}$ coefficient bounds preserves this conclusion.

By \eqref{pk:screen-lower}, on the cosine-good set the leading
norm is at least $\tfrac14\amp\lambda^2e^{-o(T)}$ uniformly
throughout $\mathcal Q_T$. Its ratio to the linear lower-order
error tends to infinity, since
$\lambda e^{-o(T)}\to\infty$, and its ratio to $\amp$ also
tends to infinity. This proves \eqref{at:signal}.
The argument compares two curvature tensors at common numerical
coordinate positions; the background curvature is never presumed
small or bounded independently of $T$.
\end{proof}

\subsection{Canonical terminal states and finite persistence}

Let $O\in\mathcal B$ be the fixed test ball from \Cref{def:exit-test},
with $\overline O\subset\mathcal N$ for the chosen observation neighborhood.
Choose a closed concentric subball $O_0\Subset O$ such that
\begin{equation}\label{at:trim}
 \mu_d(O_0)>.95\mu_d(O).
\end{equation}
Such a choice follows by monotone convergence of the smooth
positive density. Choose a smaller closed interval
$J_0\Subset J$ of positive length. For every late level $v\ge v_0$
reached by the tube, let $\mathcal H_v$ be the background hitting map
of $O_0$, where $v_0$ is the fixed late starting level. It records
the canonical state
\begin{equation}\label{at:terminal-state}
 z=(u,\theta,\pi_u,P_1,P_2).
\end{equation}
The original seed-to-covector identification uses $w\mapsto
h(w,\cdot)$ on $\Sigma$. Thus $\mu_d$ is exactly the canonical
section measure, not a position-only measure.

\begin{proposition}[Matched seeds and finite observations]
\label{prop:matched-seeds}
For $v\in T+J_0$, use the same numerical states \eqref{at:terminal-state} in
$g_T$ on its future unit shell. Their backward geodesics give
maps $\Psi_{T,v}:O_0\to O$ with
\begin{equation}\label{at:seed-closeness}
 \sup_{v\in T+J_0}\|\Psi_{T,v}-\operatorname{incl}\|_{C^1}
                                  \longrightarrow0.
\end{equation}
Write $\delta_T$ for the left side of \eqref{at:seed-closeness}.
Each map is injective and preserves the corresponding section
measure: $\Psi_{T,v}^*\mu_{d_T}=\mu_d$ on $O_0$. In particular
$\mathcal H_v(O_0)$ has canonical volume $\mu_d(O_0)$ for every such $v$.

Put
\begin{equation}\label{at:frame}
 r_T=e^{-\kappa T},\qquad
 F=(\partial_u,r_T^{-1}\partial_v,
                    \partial_{\theta^1},\partial_{\theta^2}).
\end{equation}
For $g$ and $g_T$ on the larger observation box, the Gram
matrices in $F$ and their inverses are uniformly bounded and
nondegenerate. Their ordinary coordinate jets through order
thirty are $e^{o(T)}$. The selected unit tangent is bounded
in $F$, all state momenta lie in one bounded box, and
\begin{equation}\label{at:volume}
 c r_T\,\dd u\,\dd v\,\dd^2\theta
 \le \dd\vol_g,\ \dd\vol_{g_T}
 \le C r_T\,\dd u\,\dd v\,\dd^2\theta.
\end{equation}

For each fixed finite $T$ and each $\eta_{\rm map}>0$, these observations
persist for all data in a sufficiently small relative weighted smooth
neighborhood of $d_T$, with signal lower bound $\amp$ in place of
$2\amp$. The persisted common coordinates need not be optical.
The canonical section volume, injective maps into $O$, and frame
and density bounds persist. Their backward maps satisfy
\[
 \sup_{v\in T+J_0}\|\widehat\Psi_{T,v}-\Psi_{T,v}\|_{C^1}
 <\eta_{\rm map},
 \qquad
 \sup_{v\in T+J_0}\|\widehat\Psi_{T,v}-\operatorname{incl}\|_{C^1}
 \le\delta_T+\eta_{\rm map}.
\]
The neighborhood may depend on $T$ and $\eta_{\rm map}$.
\end{proposition}

\begin{proof}
The unit mass shell in double-null variables reads
\[
 p=-(\pi_u+b^AP_A)>0,\qquad
 2p\pi_v+a(1+\gamma^{AB}P_AP_B)=0.
\]
Thus the reduced $v$-Hamiltonian is
\begin{equation}\label{at:reduced-H}
 \mathcal H=-\pi_v=
        \frac{a(1+\gamma^{AB}P_AP_B)}{2p}.
\end{equation}
The same expression with primed fields defines the new future
unit root at the identical numerical terminal state. The
background tube keeps $p$ in a compact positive interval and
the other momenta bounded. Since $b_T-b\to0$ uniformly, the
new $p_T=-(\pi_u+b_T^AP_A)$ remains bounded away from zero.
Choosing the root with $\dot v=p_T/a_T>0$ fixes its future
orientation.

Write $V$ for the background Hamiltonian vector field. On the common
late state tube this field and its first two state derivatives are bounded by
$a e^{o(v)}$, by \Cref{lem:bg-branch-tube}. This bound is
integrable over the late $v$ interval. Differentiating
\eqref{at:reduced-H} shows that the new field and its first
two state derivatives differ by an absolutely exponentially
small quantity, uniformly on that tube. Indeed $a_T=q\mathcal A_T$
retains its factor $q$ under ordinary differentiation;
the denominators are powers of $p_T$ separated from zero;
and all relative coefficient differences and their needed
jets are exponentially small by \Cref{prop:packet-preparation}
and \Cref{thm:dn-comparison}. Fixed higher background factors
cost $e^{o(T)}$, which is absorbed by that exponential.

The terminal state difference is exactly zero. Backward Gronwall
on a late interval of length $O(T)$ therefore gives state
differences tending to zero. For first variations the equation
is $J'=D_zV\,J$. Its coefficient difference is bounded by the
small difference of vector-field derivatives plus
$|D_z^2V|$ times the already small state difference. The latter
background bound is integrable. A second Gronwall estimate
therefore gives uniform $C^1$ convergence of the backward maps.
The background forward hitting map also has a uniform derivative bound:
if $v_0$ is the fixed late starting level, then
\[
 \|D_w\mathcal H_v\|
 \le \|D_w\mathcal H_{v_0}\|
       \exp\!\left(\int_{v_0}^{v}\|D_zV(s)\|\,\dd s\right)\le C.
\]
Composing the terminal-state estimates with $\mathcal H_v$ therefore
preserves their $C^1$ convergence in the original seed variables.
The strict tube and construction margins prevent exit during
this comparison. Below any fixed late $v$ the remaining flow
tubes to $\mathcal S$ and then $\Sigma$ are compact and
transverse. Smooth finite-flow dependence handles their hitting
maps; protected agreement makes their metrics exactly the
background metrics for sufficiently large $T$. This proves
\eqref{at:seed-closeness} and, by the margin $O_0\Subset O$,
places the new seeds in $O$.

Here is the exact volume statement. On the cotangent bundle,
let
\[
 \omega=\dd\pi_u\wedge\dd u+\dd\pi_v\wedge\dd v
                              +\dd P_A\wedge\dd\theta^A.
\]
On the unit Hamiltonian shell its restriction has the geodesic
flow direction as kernel. A single-valued transverse hitting
map therefore preserves the restrictions of $\omega$ to two
sections: flowing by a variable time only adds multiples of its
kernel to tangent vectors. On $v=\mathrm{constant}$ this
restriction is exactly
\[
 \omega_v=\dd\pi_u\wedge\dd u+\dd P_A\wedge\dd\theta^A,
\]
independently of the shell root $\pi_v$. Its volume is
$\dd u\,\dd^2\theta\,\dd\pi_u\,\dd^2P$.
Apply this to the background hit and to the new backward hit.
It gives $\Psi_{T,v}^*\mu_{d_T}=\mu_d$ exactly.
The hits are single-valued on these compact transverse flow
tubes; uniqueness also gives injectivity. Alternatively, in
the ball coordinates \eqref{at:seed-closeness} makes the map
$\Id+f$ with $\|Df\|<1/2$, which is injective on the convex
closed ball. Nothing in this argument assumes injectivity of
the projection from seed states to observation positions.

At an observation state the unit tangent satisfies
\begin{equation}\label{at:tangent-components}
 \dot u=\frac{1+\gamma^{AB}P_AP_B}{2p},\qquad
 r_T\dot v=\frac{r_Tp}{a},\qquad
 \dot\theta=b\dot u+\gamma^{-1}P,
\end{equation}
and likewise for $g_T$. On the whole observation box, $u$ and
$v-T$ are bounded, so
$a/r_T=q_0\mathcal A e^{-\kappa(u+v-T)}$ is bounded above
and below. These formulas prove the tangent bounds in $F$.
The Gram matrix has the block form
\begin{equation}\label{at:gram}
 G_F=\begin{pmatrix}
       \gamma(b,b)&-a/r_T&-b^C\gamma_{CB}\\
       -a/r_T&0&0\\
       -\gamma_{AC}b^C&0&\gamma_{AB}
      \end{pmatrix},\qquad
 \det G_F=-(a/r_T)^2\det\gamma.
\end{equation}
The low field bounds give bounded entries and a determinant
bounded away from zero, hence a bounded inverse. The same
holds for $g_T$. All ordinary derivatives of $a/r_T$ retain
this bounded relative factor, and the other required field
jets are $e^{o(T)}$ by the individual estimates. This proves
the thirty-jet assertion on the larger box, not just on the
hitting positions. Finally
$\sqrt{|\det g|}=a\sqrt{\det\gamma}$ proves
\eqref{at:volume}.

We finish with the finite persistence assertion. For a fixed $T$
the closures of the observation box, a small enlargement, and
all backward flow tubes just used are compact subsets of the
smooth attached development with an interior margin. The
compact-subset Cauchy stability and geometric uniqueness of
\cite[\CompanionRef{foundation}{cat:closed-tests}]{KerrCTwo}, also used in the closed-test proof
of \cite[\CompanionRef{signed}{test:closed}, proof]{KerrCOne}, supply corresponding compact
coordinate witnesses for all sufficiently close smooth data.
Require closeness in the finite $C^k$ norms needed here. Since
$T$ is fixed, it can be arbitrarily accurate even in the
rescaled frame $F$. Preserve the Gram and inverse bounds,
the thirty derivative bounds (for example allowing additive
error at most one), the density comparison, and the curvature
margin, so that
\begin{equation}\label{at:persisted-signal}
 |D(\widehat g_T)-D(g)|>\amp
        \quad\hbox{when}\quad |\cos(\lambda v)|\ge1/2.
\end{equation}

The common coordinates for $\widehat g_T$ need not retain
double-null form. To define its terminal state solve
$\tfrac12\widehat g_T^{-1}(\pi,\pi)=-\tfrac12$ for
$\pi_v$ with the other six numbers fixed as in
\eqref{at:terminal-state}. At the original terminal state the
derivative with respect to $\pi_v$ is $U^v>0$.
On the fixed compact state set it is separated from zero, so
the implicit function theorem supplies the same smooth simple
future root for close metrics. The level-$v$ section remains
transverse to this flow. Smooth backward-flow dependence lets us impose
$\sup_{v\in T+J_0}\|\widehat\Psi_{T,v}-\Psi_{T,v}\|_{C^1}
<\min\{\eta_{\rm map},\eta_0\}$, where a fixed $\eta_0>0$ preserves
injectivity and the image margin in $O$.
The restriction of the canonical symplectic form to that section
is still $\omega_v$, regardless of optical gauge. Its transported
volume is therefore exactly the new initial seed measure.
All statements hold for this fixed experiment, with no assertion
of a data-neighborhood radius uniform in $T$.
\end{proof}
